\chapter{Fonts and text}

\lstinline|DrawText| with the built-in font gets you through a prototype. The
moment you want a game that looks designed, you need the rest of rtext.

\section{Loading a font properly}

\begin{lstlisting}
Font font = LoadFont("kenney.ttf");                     // default size 32
Font font = LoadFontEx("kenney.ttf", 48, NULL, 0);      // the one you want
\end{lstlisting}

\lstinline|LoadFontEx| takes the size to rasterise at, and optionally a list of
codepoints to include. Passing \lstinline|NULL, 0| gives you the default set of
95 ASCII characters.

\gotcha{A font is baked into a bitmap at \textbf{one fixed size}. Load it at 16
and draw it at 64 and it will be blurry --- that is not a bug, it is a small
image being stretched. Load it at, or slightly above, the largest size you will
actually draw.}

For accented characters, and for anything beyond English, you must say so at load
time:

\begin{lstlisting}
int codepoints[512] = { 0 };
for (int i = 0; i < 95; i++) codepoints[i] = 32 + i;        // basic ASCII
for (int i = 0; i < 96; i++) codepoints[95 + i] = 0xC0 + i; // Latin-1 supplement

Font font = LoadFontEx("font.ttf", 32, codepoints, 191);
\end{lstlisting}

A character you did not ask for simply does not render. If your accents come out
as blanks or question marks, this is why.

\section{Drawing with it}

\begin{lstlisting}
DrawText(text, x, y, size, colour);                          // default font only
DrawTextEx(font, text, position, size, spacing, colour);     // any font
DrawTextPro(font, text, position, origin, rot, size, spacing, colour);
\end{lstlisting}

\lstinline|DrawTextEx| is the everyday one. The \lstinline|spacing| parameter is
extra pixels between characters, and it matters more than people expect: a value
of 1 or 2 makes most fonts noticeably easier to read on screen.

\lstinline|DrawTextPro| adds an origin and a rotation, for damage numbers that
float and spin, or text that follows a curve.

\section{Measuring, which is how you lay anything out}

There is no alignment parameter anywhere in raylib. You measure and subtract:

\begin{lstlisting}
int w = MeasureText(txt, 20);                                // default font
Vector2 size = MeasureTextEx(font, txt, 20, 2);              // any font, gives w AND h

DrawText(txt, GetScreenWidth()/2 - w/2, y, 20, BLACK);       // centred
DrawText(txt, GetScreenWidth() - w - 20, y, 20, BLACK);      // right aligned
\end{lstlisting}

\lstinline|MeasureTextEx| returning a \lstinline|Vector2| is what lets you draw a
background box that fits its text, which is most of what a UI is:

\begin{lstlisting}
Vector2 s = MeasureTextEx(font, label, 20, 2);
DrawRectangle(x - 8, y - 6, (int)s.x + 16, (int)s.y + 12, Fade(BLACK, 0.7f));
DrawTextEx(font, label, (Vector2){ x, y }, 20, 2, RAYWHITE);
\end{lstlisting}

\section{The string helpers}

rtext also contains a small string library, which exists because doing this in
plain C is tedious:

\begin{center}
\begin{tabular}{@{}ll@{}}
\toprule
\lstinline|TextFormat("%i/%i", a, b)| & sprintf, managed buffer \\
\lstinline|TextLength(s)| & like strlen \\
\lstinline|TextSubtext(s, pos, len)| & a slice \\
\lstinline|TextReplace(s, from, to)| & search and replace \\
\lstinline|TextInsert(s, insert, pos)| & \\
\lstinline|TextJoin(list, count, sep)| & \\
\lstinline|TextSplit(s, delim, &count)| & split into pieces \\
\lstinline|TextFindIndex(s, find)| & position, or $-1$ \\
\lstinline|TextToUpper / TextToLower| & \\
\lstinline|TextToInteger / TextToFloat| & parse \\
\lstinline|TextCopy(dst, src)| & returns the byte count \\
\bottomrule
\end{tabular}
\end{center}

\gotcha{\lstinline|TextFormat|, \lstinline|TextSubtext|, \lstinline|TextReplace|
and most of the others write into \textbf{internal rotating buffers}. The pointer
is valid for drawing right now and becomes garbage after a few more calls. Never
store one. If you need to keep the string, \lstinline|TextCopy| it into your own
array.

There is one exception that goes the other way: \lstinline|TextReplace| in some
versions returns heap memory you must free. When in doubt, check the comment in
\file{raylib.h} --- it says so per function.}

\section{Unicode, briefly}

In UTF-8 one character is not one byte. \lstinline|TextLength| counts
\emph{bytes}; the number of characters --- codepoints --- can be smaller:

\begin{lstlisting}
int count = 0;
int *codepoints = LoadCodepoints(utf8Text, &count);
// count is the number of CHARACTERS, TextLength() is the number of BYTES
UnloadCodepoints(codepoints);
\end{lstlisting}

This matters the day you let a player type their name and then try to delete one
character with backspace. Deleting one byte from a multi-byte character produces
garbage. \lstinline|GetCodepointPrevious()| gives you the correct amount to step
back.

\section{Where to get fonts}

Any \file{.ttf} or \file{.otf} works. For games, pixel fonts hold up best at
small sizes; Kenney's free asset packs and Google Fonts both have plenty that are
free to ship. Load them with \lstinline|LoadFontEx| at the exact size you will
draw, and set \lstinline|TEXTURE_FILTER_POINT| on \lstinline|font.texture| if it
is a pixel font, so it stays crisp.

\tryit{Module 7 draws the same string through all three functions, shows
measuring and centring live, runs through every string helper with its output,
and counts bytes against codepoints on a UTF-8 string.}{./cheatsheet/build.sh 7}
